% Part: proof-theory
% Chapter: normalization
% Section: permutations

\documentclass[../../../include/open-logic-section]{subfiles}

\begin{document}

\olfileid{pt}{nor}{per}

\olsection{في تحويلات التبديل}

القطع في~\Log{N1i} هو المقطع الذي ينتهي بالمقدمة الكبرى لقاعدة
!!a{elimination}. ومن عمل برهان التطبيع تقصير القطوع القصوى، لينقص بذلك
طول القطع في~!!{proof}. وإذا زاد طول القطع على~$\OLMathNumeral{1}$ تعددت صور نهايته؛
وذلك بحسب كون آخر ظهور لـ~!!{formula}~$!A$ فيه نتيجة لـ
(\Elim{\lor} أو~\Elim{\lexists})، وبحسب استدلال~!!{elimination} الذي
تكون تلك الصيغة مقدمته الكبرى.

فإذا كان الاستدلالان، مثلًا،~\Elim{\lor} و\Elim{\land}، وكانت
$!A \ident !B \land !C$، وكان~!!{proof} الجزئي~$\delta_\OLMathNumeral{1}$ منتهيًا هكذا:
\[
\AxiomC{}
\RightLabel{$\delta_\OLMathNumeral{2}$}
\DeduceC{$!D \lor !E$}
\AxiomC{$\Discharge{!D}{\OLId{latin-ordinary}{x}}$}
\RightLabel{$\delta_\OLMathNumeral{3}$}
\DeduceC{$!B \land !C$}
\AxiomC{$\Discharge{!E}{\OLId{latin-ordinary}{y}}$}
\RightLabel{$\delta_\OLMathNumeral{4}$}
\DeduceC{$!B \land !C$}
\DischargeRule{\Elim{\lor}}{\OLId{latin-ordinary}{x}\,\OLId{latin-ordinary}{y}}
\TrinaryInfC{$!B \land !C$}
\RightLabel{\Elim{\land}[\OLMathNumeral{1}]}
\UnaryInfC{$!B$}
\DisplayProof
\]
كان آخر ظهور لـ~!!{formula}~$!A_\OLId{latin-ordinary}{n}$ في القطع نتيجة~\Elim{\lor}، وهو
الظهور الأسفل لـ~$!B \land !C$. وأما الظهور السابق له من~!!{formula}
$!A_{\OLId{latin-ordinary}{n}-\OLMathNumeral{1}}$ فهو إحدى المقدمتين الصغريين لـ~\Elim{\lor}.

ولتقصير القطع نعكس ترتيب استدلالي~\Elim{\lor} و\Elim{\land}، أي
«نبدلهما». فنقيم مقام البرهان الجزئي~$\delta_\OLMathNumeral{1}$، داخل~$\delta$، ما يأتي:
\[
\AxiomC{}
\RightLabel{$\delta_\OLMathNumeral{2}$}
\DeduceC{$!D \lor !E$}
\AxiomC{$\Discharge{!D}{\OLId{latin-ordinary}{x}}$}
\RightLabel{$\delta_\OLMathNumeral{3}$}
\DeduceC{$!B \land !C$}
\RightLabel{\Elim{\land}[\OLMathNumeral{1}]}
\UnaryInfC{$!B$}
\AxiomC{$\Discharge{!E}{\OLId{latin-ordinary}{y}}$}
\RightLabel{$\delta_\OLMathNumeral{4}$}
\DeduceC{$!B \land !C$}
\RightLabel{\Elim{\land}[\OLMathNumeral{1}]}
\UnaryInfC{$!B$}
\DischargeRule{\Elim{\lor}}{\OLId{latin-ordinary}{x}\,\OLId{latin-ordinary}{y}}
\TrinaryInfC{$!B$}
\DisplayProof
\]
فالقطوع التي كانت نهايتها مقدمة~\Elim{\land} تحت~\Elim{\lor} صارت
نهايتها مقدمة أحد استدلالي~\Elim{\land} الواقعين فوق المقدمتين الصغريين
لـ~\Elim{\lor}؛ فنقص طول كل منها بمقدار~$\OLMathNumeral{1}$.

ثم إن كل قطع واقع كله في~$\delta_\OLMathNumeral{2}$ أو~$\delta_\OLMathNumeral{3}$ أو~$\delta_\OLMathNumeral{4}$، أو
واقع كله خارج~$\delta_\OLMathNumeral{1}$ وداخل~$\delta$، باق على حاله: يشتمل على
!!{formula}s نفسها ويمر بالاستدلالات نفسها؛ فرتبته وطوله هما رتبة القطع
المقابل في~$\delta$ وطوله. فلم تزد رتبة القطع في~!!{proof} الجديد، وصار
طول القطع في~!!{proof} الجديد دون طوله في~!!{proof} القديم.

\begin{prob}
ينقص طول قطع واحد على الأقل بمقدار~$\OLMathNumeral{1}$ بعد تحويل تبديل. وإذا بدّلنا
\Elim{\land} عبر~\Elim{\lor}، نقص، في الحقيقة، طول مقطعي قطع، ومن ثم
ينقص طول القطع في~!!{proof} بمقدار~$\OLMathNumeral{2}$ على الأقل. فهل يجوز أن ينقص طول
القطع في~!!a{proof} بأكثر من~$\OLMathNumeral{2}$ بتحويل واحد من هذا النوع؟ وهل يجوز أن
تنقص \emph{رتبة} القطع؟
\end{prob}

ونسمي نقل~!!a{elimination} إلى ما فوق قاعدة~\Elim{\lor} أو
\Elim{\lexists} \emph{تحويل تبديل}. وقد أثبتت أنماط بعض أزواج القواعد
في~\olref{tab:permutation-conversions}، وتركت بقيتها تمارين.

% \begin{table}
% \begin{multline*}
% \shoveleft{\AxiomC{$!D \lor !E$}
% \AxiomC{$\Discharge{!D}{x}$}
% \shortDeduce
% \DeduceC{$!B \land !C$}
% \AxiomC{$\Discharge{!E}{y}$}
% \shortDeduce
% \DeduceC{$!B \land !C$}
% \DischargeRule{\Elim{\lor}}{x\,y}
% \TrinaryInfC{$!B \land !C$}
% \RightLabel{\Elim{\land}[1]}
% \UnaryInfC{$!B$}
% \DisplayProof}\\
% \shoveright{\longrightarrow\ 
% \AxiomC{$!D \lor !E$}
% \AxiomC{$\Discharge{!D}{x}$}
% \shortDeduce
% \DeduceC{$!B \land !C$}
% \RightLabel{\Elim{\land}[1]}
% \UnaryInfC{$!B$}
% \AxiomC{$\Discharge{!E}{y}$}
% \shortDeduce
% \DeduceC{$!B \land !C$}
% \RightLabel{\Elim{\land}[1]}
% \UnaryInfC{$!B$}
% \DischargeRule{\Elim{\lor}}{x\,y}
% \TrinaryInfC{$!B$}
% \DisplayProof}
% \\
% \shoveleft{\AxiomC{$!D \lor !E$}
% \AxiomC{$\Discharge{!D}{x}$}
% \shortDeduce
% \DeduceC{$!B \lif !C$}
% \AxiomC{$\Discharge{!E}{y}$}
% \shortDeduce
% \DeduceC{$!B \lif !C$}
% \DischargeRule{\Elim{\lor}}{x\,y}
% \TrinaryInfC{$!B \lif !C$}
% \AxiomC{$!B$}
% \RightLabel{\Elim{\lif}}
% \BinaryInfC{$!C$}
% \DisplayProof}\\
% \shoveright{\longrightarrow\ 
% \AxiomC{$!D \lor !E$}
% \AxiomC{$\Discharge{!D}{x}$}
% \shortDeduce
% \DeduceC{$!B \lif !C$}
% \AxiomC{$!B$}
% \RightLabel{\Elim{\lif}}
% \BinaryInfC{$!C$}
% \AxiomC{$\Discharge{!E}{y}$}
% \shortDeduce
% \DeduceC{$!B \lif !C$}
% \AxiomC{$!B$}
% \RightLabel{\Elim{\lif}}
% \BinaryInfC{$!C$}
% \DischargeRule{\Elim{\lor}}{x\,y}
% \TrinaryInfC{$!C$}
% \DisplayProof}
% \\
% \shoveleft{\AxiomC{$!D \lor !E$}
% \AxiomC{$\Discharge{!D}{x}$}
% \shortDeduce
% \DeduceC{$!B \lif !C$}
% \AxiomC{$\Discharge{!E}{y}$}
% \shortDeduce
% \DeduceC{$!B \lif !C$}
% \DischargeRule{\Elim{\lor}}{x\,y}
% \TrinaryInfC{$!B \lif !C$}
% \AxiomC{$!B$}
% \RightLabel{\Elim{\lif}}
% \BinaryInfC{$!C$}
% \DisplayProof}\\
% \shoveright{\longrightarrow\ 
% \AxiomC{$!D \lor !E$}
% \AxiomC{$\Discharge{!D}{x}$}
% \shortDeduce
% \DeduceC{$!B \lif !C$}
% \AxiomC{$!B$}
% \RightLabel{\Elim{\lif}}
% \BinaryInfC{$!C$}
% \AxiomC{$\Discharge{!E}{y}$}
% \shortDeduce
% \DeduceC{$!B \lif !C$}
% \AxiomC{$!B$}
% \RightLabel{\Elim{\lif}}
% \BinaryInfC{$!C$}
% \DischargeRule{\Elim{\lor}}{x\,y}
% \TrinaryInfC{$!C$}
% \DisplayProof}
% \\
% \shoveleft{\AxiomC{$!D \lor !E$}
% \AxiomC{$\Discharge{!D}{x}$}
% \shortDeduce
% \DeduceC{$\lforall[x][!B(x)]$}
% \AxiomC{$\Discharge{!E}{y}$}
% \shortDeduce
% \DeduceC{$\lforall[x][!B(x)]$}
% \DischargeRule{\Elim{\lor}}{x\,y}
% \TrinaryInfC{$\lforall[x][!B(x)]$}
% \RightLabel{\Elim{\lforall}}
% \UnaryInfC{$!B(t)$}
% \DisplayProof}\\
% \shoveright{\longrightarrow\ 
% \AxiomC{$!D \lor !E$}
% \AxiomC{$\Discharge{!D}{x}$}
% \shortDeduce
% \DeduceC{$\lforall[x][!B(x)]$}
% \RightLabel{\Elim{\lforall}}
% \UnaryInfC{$!B(t)$}
% \AxiomC{$\Discharge{!E}{y}$}
% \shortDeduce
% \DeduceC{$\lforall[x][!B(x)]$}
% \RightLabel{\Elim{\lforall}}
% \UnaryInfC{$!B(t)$}
% \DischargeRule{\Elim{\lor}}{x\,y}
% \TrinaryInfC{$!B(t)$}
% \DisplayProof}
% \end{multline*}
% \caption{Some permutation conversions for \Log{N1i}.}
% \ollabel{tab:permutation-conversions}
% \end{table}

{\small
\begin{longtable}{@{}l@{}}
\AxiomC{$!D \lor !E$}
\AxiomC{$\Discharge{!D}{\OLId{latin-ordinary}{x}}$}
\shortDeduce
\DeduceC{$!B \land !C$}
\AxiomC{$\Discharge{!E}{\OLId{latin-ordinary}{y}}$}
\shortDeduce
\DeduceC{$!B \land !C$}
\DischargeRule{\Elim{\lor}}{x\,y}
\TrinaryInfC{$!B \land !C$}
\RightLabel{\Elim{\land}[1]}
\UnaryInfC{$!B$}
\DisplayProof
$\longrightarrow$\\[6ex]
\quad\AxiomC{$!D \lor !E$}
\AxiomC{$\Discharge{!D}{\OLId{latin-ordinary}{x}}$}
\shortDeduce
\DeduceC{$!B \land !C$}
\RightLabel{\Elim{\land}[1]}
\UnaryInfC{$!B$}
\AxiomC{$\Discharge{!E}{\OLId{latin-ordinary}{y}}$}
\shortDeduce
\DeduceC{$!B \land !C$}
\RightLabel{\Elim{\land}[1]}
\UnaryInfC{$!B$}
\DischargeRule{\Elim{\lor}}{x\,y}
\TrinaryInfC{$!B$}
\DisplayProof
\\[10ex]
\AxiomC{$!D \lor !E$}
\AxiomC{$\Discharge{!D}{\OLId{latin-ordinary}{x}}$}
\shortDeduce
\DeduceC{$!B \lif !C$}
\AxiomC{$\Discharge{!E}{\OLId{latin-ordinary}{y}}$}
\shortDeduce
\DeduceC{$!B \lif !C$}
\DischargeRule{\Elim{\lor}}{x\,y}
\TrinaryInfC{$!B \lif !C$}
\AxiomC{$!B$}
\RightLabel{\Elim{\lif}}
\BinaryInfC{$!C$}
\DisplayProof
$\longrightarrow$\\[6ex]
\AxiomC{$!D \lor !E$}
\AxiomC{$\Discharge{!D}{\OLId{latin-ordinary}{x}}$}
\shortDeduce
\DeduceC{$!B \lif !C$}
\AxiomC{$!B$}
\RightLabel{\Elim{\lif}}
\BinaryInfC{$!C$}
\AxiomC{$\Discharge{!E}{\OLId{latin-ordinary}{y}}$}
\shortDeduce
\DeduceC{$!B \lif !C$}
\AxiomC{$!B$}
\RightLabel{\Elim{\lif}}
\BinaryInfC{$!C$}
\DischargeRule{\Elim{\lor}}{x\,y}
\TrinaryInfC{$!C$}
\DisplayProof
\\[10ex]
\AxiomC{$!D \lor !E$}
\AxiomC{$\Discharge{!D}{\OLId{latin-ordinary}{x}}$}
\shortDeduce
\DeduceC{$\lforall[\OLId{latin-ordinary}{x}][!B(\OLId{latin-ordinary}{x})]$}
\AxiomC{$\Discharge{!E}{\OLId{latin-ordinary}{y}}$}
\shortDeduce
\DeduceC{$\lforall[\OLId{latin-ordinary}{x}][!B(\OLId{latin-ordinary}{x})]$}
\DischargeRule{\Elim{\lor}}{x\,y}
\TrinaryInfC{$\lforall[\OLId{latin-ordinary}{x}][!B(\OLId{latin-ordinary}{x})]$}
\RightLabel{\Elim{\lforall}}
\UnaryInfC{$!B(\OLId{latin-ordinary}{t})$}
\DisplayProof
\quad$\longrightarrow$\\[8ex]
\quad\AxiomC{$!D \lor !E$}
\AxiomC{$\Discharge{!D}{\OLId{latin-ordinary}{x}}$}
\shortDeduce
\DeduceC{$\lforall[\OLId{latin-ordinary}{x}][!B(\OLId{latin-ordinary}{x})]$}
\RightLabel{\Elim{\lforall}}
\UnaryInfC{$!B(\OLId{latin-ordinary}{t})$}
\AxiomC{$\Discharge{!E}{\OLId{latin-ordinary}{y}}$}
\shortDeduce
\DeduceC{$\lforall[\OLId{latin-ordinary}{x}][!B(\OLId{latin-ordinary}{x})]$}
\RightLabel{\Elim{\lforall}}
\UnaryInfC{$!B(\OLId{latin-ordinary}{t})$}
\DischargeRule{\Elim{\lor}}{x\,y}
\TrinaryInfC{$!B(\OLId{latin-ordinary}{t})$}
\DisplayProof
\\
\caption{بعض تحويلات التبديل في~\Log{N1i}.}
\ollabel{tab:permutation-conversions}
\end{longtable}
}

\begin{prob}
اكتب تحويلات التبديل لـ~\Elim{\lor} متبوعةً بـ~\Elim{\lor} أو
\Elim{\lnot}.
\end{prob}

\begin{prob}
اكتب تحويلات التبديل لـ~\Elim{\lexists} متبوعةً بـ~\Elim{\lexists}.
واشرح لماذا لا يُنتهك شرط المتغير الذاتي في الحالة الثانية.
\end{prob}

ويبقى أن نتحقق من بقاء شرط المتغير الذاتي بعد تحويل التبديل. فلننظر في:
\[
\AxiomC{}
\RightLabel{$\delta_\OLMathNumeral{2}$}
\DeduceC{$!D \lor !E$}
\AxiomC{$\Discharge{!D}{\OLId{latin-ordinary}{x}}$}
\RightLabel{$\delta_\OLMathNumeral{3}$}
\DeduceC{$\lexists[\OLId{latin-ordinary}{x}][!B(\OLId{latin-ordinary}{x})]$}
\AxiomC{$\Discharge{!E}{\OLId{latin-ordinary}{y}}$}
\RightLabel{$\delta_\OLMathNumeral{4}$}
\DeduceC{$\lexists[\OLId{latin-ordinary}{x}][!B(\OLId{latin-ordinary}{x})]$}
\DischargeRule{\Elim{\lor}}{\OLId{latin-ordinary}{x}\,\OLId{latin-ordinary}{y}}
\TrinaryInfC{$\lexists[\OLId{latin-ordinary}{x}][!B(\OLId{latin-ordinary}{x})]$}
\AxiomC{$\Discharge{!B(\OLId{latin-ordinary}{c})}{\OLId{latin-ordinary}{z}}$}
\RightLabel{$\delta_\OLMathNumeral{5}$}
\DeduceC{$!C$}
\DischargeRule{\Elim{\lexists}}{\OLId{latin-ordinary}{z}}
\BinaryInfC{$!C$}
\DisplayProof
\]
وقبل أن نستبدل بهذا البرهان نأخذ من~$\delta_\OLMathNumeral{5}$ نسختين، ونعيد تسمية
المتغيرات الذاتية ووسوم التفريغ فيهما اتقاء الأسر، حتى يكون الزوجان
$(\OLId{latin-ordinary}{c}_\OLId{latin-looped}{D},\OLId{latin-ordinary}{z}_\OLId{latin-looped}{D})$ و$(\OLId{latin-ordinary}{c}_\OLId{latin-looped}{E},\OLId{latin-ordinary}{z}_\OLId{latin-looped}{E})$ حديثين ومتميزين زوجيًا. ثم نقيم مكانه ما يأتي:
\[
\AxiomC{}
\RightLabel{$\delta_\OLMathNumeral{2}$}
\DeduceC{$!D \lor !E$}
\AxiomC{$\Discharge{!D}{\OLId{latin-ordinary}{x}}$}
\RightLabel{$\delta_\OLMathNumeral{3}$}
\DeduceC{$\lexists[\OLId{latin-ordinary}{x}][!B(\OLId{latin-ordinary}{x})]$}
\AxiomC{$\Discharge{!B(\OLId{latin-ordinary}{c}_\OLId{latin-looped}{D})}{\OLId{latin-ordinary}{z}_\OLId{latin-looped}{D}}$}
\RightLabel{$\delta_\OLMathNumeral{5}^{(\OLId{latin-looped}{D})}$}
\DeduceC{$!C$}
\DischargeRule{\Elim{\lexists}}{\OLId{latin-ordinary}{z}_\OLId{latin-looped}{D}}
\BinaryInfC{$!C$}
\AxiomC{$\Discharge{!E}{\OLId{latin-ordinary}{y}}$}
\RightLabel{$\delta_\OLMathNumeral{4}$}
\DeduceC{$\lexists[\OLId{latin-ordinary}{x}][!B(\OLId{latin-ordinary}{x})]$}
\AxiomC{$\Discharge{!B(\OLId{latin-ordinary}{c}_\OLId{latin-looped}{E})}{\OLId{latin-ordinary}{z}_\OLId{latin-looped}{E}}$}
\RightLabel{$\delta_\OLMathNumeral{5}^{(\OLId{latin-looped}{E})}$}
\DeduceC{$!C$}
\DischargeRule{\Elim{\lexists}}{\OLId{latin-ordinary}{z}_\OLId{latin-looped}{E}}
\BinaryInfC{$!C$}
\DischargeRule{\Elim{\lor}}{\OLId{latin-ordinary}{x}\,\OLId{latin-ordinary}{y}}
\TrinaryInfC{$!C$}
\DisplayProof
\]
وقد يخشى أن يظهر المتغير الذاتي~$\OLId{latin-ordinary}{c}$، مثلًا، في~$!D$. والفرض~$!D$ مفرغ
فوق~\Elim{\lexists} الأصلي، فلا يقع في~!!{proof} الأصلي فرضًا مفتوحًا
فوق المقدمة الكبرى لاستدلال~\Elim{\lexists}، ومن ثم يصح شرط المتغير
الذاتي هناك. وأما في~!!{proof} الجديد فلا يفرغ~$!D$ إلا بعد استدلالي
\Elim{\lexists}، فيلوح أن الشرط قد اختل. غير أنا نفرض انتظام~!!{proof}s:
فكل متغير ذاتي مختص باستدلال واحد، ولا يظهر في~!!{proof} إلا فوق المقدمة
الصغرى لـ~\Elim{\lexists}. وعلى الخصوص لا يظهر~$\OLId{latin-ordinary}{c}$ في~$\delta_\OLMathNumeral{3}$ ولا
$\delta_\OLMathNumeral{4}$؛ ثم إن إعادة التسمية بالمتغيرات الحديثة تمنع الأسر بين
النسختين.

وفي المثال فائدة أخرى: فإن في~!!{proof} الجديد نسختين من~$\delta_\OLMathNumeral{5}$.
فإذا اشتمل~$\delta_\OLMathNumeral{5}$ على قطع أقصى، لم يلزم نقصان طول القطع، بل ربما زاد
بالنسخ. فلا نعمل تحويل التبديل إلا حيث نعلم خلو~$\delta_\OLMathNumeral{5}$ من قطع أقصى.
ونضمن هذا بأن نختار دائمًا قطعًا أقصى \emph{علويًا}؛ أي قطعًا لا يشتمل
!!{proof}ه الجزئي، وهو~$\delta_\OLMathNumeral{1}$ هنا، على قطع أقصى ينتهي فوق آخر استدلال
في ذلك~!!{proof} الجزئي. وبذلك يمتنع وجود قطع أقصى آخر في~$\delta_\OLMathNumeral{5}$، بل
في~$\delta_\OLMathNumeral{2}$ أو~$\delta_\OLMathNumeral{3}$ أو~$\delta_\OLMathNumeral{4}$ أيضًا.

\end{document}
